\documentclass[a4paper,12pt]{article}
\usepackage[utf8]{inputenc}
\usepackage[T1]{fontenc}
\usepackage[french]{babel}
\usepackage{geometry}
\usepackage{booktabs}
\usepackage{array}
\usepackage{xcolor}
\usepackage{enumitem}
\usepackage{titlesec}
\usepackage{hyperref}
\usepackage{longtable}
\usepackage{parskip}
\usepackage{tikz}
\usepackage{pgf}
\usepackage{fancyhdr}
\usepackage{mdframed}
\usepackage{tabularx}
\usepackage{multirow}
\usepackage{colortbl}

\usetikzlibrary{shapes.geometric, arrows.meta, positioning, fit, backgrounds, shadows}

\geometry{margin=2.5cm}

% ── Colors ───────────────────────────────────────────────────────────
\definecolor{primary}{RGB}{99,77,189}
\definecolor{accent}{RGB}{60,40,130}
\definecolor{bronze}{RGB}{176,120,60}
\definecolor{silver}{RGB}{120,130,150}
\definecolor{gold}{RGB}{195,155,30}
\definecolor{opmcolor}{RGB}{190,50,50}
\definecolor{pmacolor}{RGB}{50,130,190}
\definecolor{ptecolor}{RGB}{40,160,100}
\definecolor{execcolor}{RGB}{140,60,180}
\definecolor{factbg}{RGB}{240,245,255}
\definecolor{dimbg}{RGB}{255,248,235}
\definecolor{confbg}{RGB}{235,255,240}
\definecolor{lightpurple}{RGB}{245,242,255}

% ── Section style ────────────────────────────────────────────────────
\titleformat{\section}{\large\bfseries\color{accent}}{}{0em}{}[\titlerule]
\titleformat{\subsection}{\normalsize\bfseries\color{primary}}{}{0em}{\quad}
\titleformat{\subsubsection}{\small\bfseries\color{primary}}{}{0em}{\qquad}

\hypersetup{colorlinks=true, linkcolor=primary, urlcolor=primary}

% ── Header / Footer ──────────────────────────────────────────────────
\pagestyle{fancy}
\fancyhf{}
\rhead{\small\color{primary} Architecture ETL \& Data Warehouse}
\lhead{\small\color{primary} Reporting Hub}
\rfoot{\small\thepage}
\renewcommand{\headrulewidth}{0.4pt}

% ── TikZ styles ──────────────────────────────────────────────────────
\tikzset{
  factbox/.style={
    rectangle, rounded corners=4pt,
    fill=factbg, draw=primary, line width=1.5pt,
    text width=3.2cm, align=center, font=\small\bfseries,
    drop shadow
  },
  dimbox/.style={
    rectangle, rounded corners=3pt,
    fill=dimbg, draw=bronze!70, line width=1pt,
    text width=2.8cm, align=center, font=\footnotesize
  },
  confdim/.style={
    rectangle, rounded corners=3pt,
    fill=confbg, draw=ptecolor!80, line width=1.5pt,
    text width=2.8cm, align=center, font=\footnotesize\bfseries
  },
  arrow/.style={-Latex, draw=gray!60, line width=0.8pt},
  confarrow/.style={-Latex, draw=ptecolor!70, line width=1pt, dashed},
  layerbox/.style={
    rectangle, rounded corners=6pt,
    draw=#1!60, fill=#1!10, line width=1.5pt,
    text width=2.6cm, minimum height=1.4cm, align=center,
    font=\small\bfseries, #1
  }
}

% ─────────────────────────────────────────────────────────────────────
\title{
  \vspace{-1cm}
  {\color{primary}\rule{\linewidth}{2pt}}\\[0.4em]
  \textbf{Reporting Hub}\\
  {\large Architecture Technique, Pipeline ETL\\[0.2em]
         \& Modélisation Data Warehouse}\\[0.3em]
  {\color{primary}\rule{\linewidth}{1pt}}
}
\author{Hamza Baktache}
\date{Juillet 2026}

\begin{document}
\maketitle
\tableofcontents
\newpage

% ═════════════════════════════════════════════════════════════════════
\section{Stack Technologique}

\begin{center}
\begin{tabular}{>{\bfseries\color{primary}}p{3.5cm} p{4cm} p{6.5cm}}
\toprule
\textbf{Couche} & \textbf{Technologie} & \textbf{Rôle} \\
\midrule
\multicolumn{3}{l}{\cellcolor{lightpurple}\small\bfseries \quad Sources de données} \\
\quad Données brutes & Kaggle (datasets publics)  & Jeux de données réels importés dans les trois backends pour simuler un environnement d'entreprise \\
\quad OPM Backend   & Node.js + Express + MongoDB & API REST du système de gestion des tickets/opérations \\
\quad PTE Backend   & Node.js + Express + MongoDB & API REST du système de gestion RH \& terrain \\
\quad PMA Backend   & Node.js + Express + MongoDB & API REST du système de gestion de projets \\
\midrule
\multicolumn{3}{l}{\cellcolor{lightpurple}\small\bfseries \quad Data Warehouse \& ETL} \\
\quad ETL Engine    & Python 3.11 + FastAPI       & Extraction, transformation, chargement (pipeline Medallion) \\
\quad Base Argent   & MongoDB 7                   & Couche Bronze (données brutes) et Argent (données normalisées) \\
\quad Base Or       & PostgreSQL 16               & Dimensions, tables de faits et vues matérialisées (data marts) \\
\quad Scheduler     & APScheduler                 & Planification automatique des runs ETL \\
\midrule
\multicolumn{3}{l}{\cellcolor{lightpurple}\small\bfseries \quad API \& Sécurité} \\
\quad API Reporting & FastAPI 0.110 + Python-JOSE  & Endpoints KPI, rapports, auth JWT (Bearer Token) \\
\quad Auth          & bcrypt + JWT (HS256)         & Hachage des mots de passe, tokens d'accès \\
\midrule
\multicolumn{3}{l}{\cellcolor{lightpurple}\small\bfseries \quad Frontend} \\
\quad UI            & Angular 17 (Standalone)     & SPA avec composants autonomes, routing lazy-load \\
\quad Graphiques    & Chart.js + ng2-charts        & Visualisation des KPI (courbes, barres, camembert) \\
\quad Styles        & Tailwind CSS                 & Utilitaires CSS, design Material You \\
\quad Réactivité    & RxJS + Angular Signals       & Flux de données réactifs et state management \\
\midrule
\multicolumn{3}{l}{\cellcolor{lightpurple}\small\bfseries \quad Infrastructure} \\
\quad Conteneurs    & Docker + Docker Compose      & Orchestration de 8 services (OPM, PTE, PMA, reporting, seeder, BDD) \\
\quad Mail          & MailHog                      & Serveur SMTP de test pour les rapports planifiés \\
\bottomrule
\end{tabular}
\end{center}

% ═════════════════════════════════════════════════════════════════════
\section{Pipeline ETL — Architecture Medallion}

L'ETL suit l'architecture \textbf{Medallion} à trois couches, garantissant la traçabilité,
la rejouabilité et la qualité progressive des données.

\vspace{1em}

\begin{center}
\begin{tikzpicture}[node distance=1.6cm, >=Latex]
  % Layers
  \node[layerbox=bronze, text=bronze] (bronze) {BRONZE\\{\footnotesize Données brutes}};
  \node[layerbox=silver, text=silver, right=2cm of bronze] (silver) {ARGENT\\{\footnotesize Normalisé}};
  \node[layerbox=gold,   text=gold,   right=2cm of silver] (gold)   {OR\\{\footnotesize Dimensionnel}};
  \node[layerbox=primary, text=primary, right=2cm of gold] (mart)   {MART\\{\footnotesize Agrégats KPI}};

  % Sources
  \node[above=1cm of bronze, xshift=-0.5cm, font=\tiny\bfseries, text=opmcolor] (opm) {OPM API};
  \node[above=1cm of bronze, font=\tiny\bfseries, text=ptecolor] (pte) {PTE API};
  \node[above=1cm of bronze, xshift=0.5cm, font=\tiny\bfseries, text=pmacolor] (pma) {PMA API};

  \draw[->, draw=opmcolor!60, line width=0.8pt] (opm) -- (bronze);
  \draw[->, draw=ptecolor!60, line width=0.8pt] (pte) -- (bronze);
  \draw[->, draw=pmacolor!60, line width=0.8pt] (pma) -- (bronze);

  % Flow
  \draw[->, line width=1.2pt, draw=bronze!80] (bronze) -- node[above, font=\tiny]{normalise} (silver);
  \draw[->, line width=1.2pt, draw=silver!80] (silver) -- node[above, font=\tiny]{MDM + dims} (gold);
  \draw[->, line width=1.2pt, draw=gold!80]   (gold)   -- node[above, font=\tiny]{vues mat.} (mart);

  % Storage labels
  \node[below=0.4cm of bronze, font=\tiny, text=bronze] {MongoDB (raw\_*)};
  \node[below=0.4cm of silver, font=\tiny, text=silver] {MongoDB (stg\_*)};
  \node[below=0.4cm of gold,   font=\tiny, text=gold]   {PostgreSQL (core.*)};
  \node[below=0.4cm of mart,   font=\tiny, text=primary]{PostgreSQL (mart\_*)};
\end{tikzpicture}
\end{center}

\vspace{0.5em}

\subsection{Couche Bronze — Extraction}

\subsubsection{Origine des données — Kaggle}

Les données métier utilisées dans ce projet proviennent de \textbf{jeux de données publics
issus de la plateforme Kaggle}. Ces datasets ont été sélectionnés pour leur représentativité
des domaines OPM, PMA et PTE, puis intégrés dans les trois applications sources (backends
Node.js) qui les exposent via des API REST simulant un environnement d'entreprise réel.
Cette approche permet de disposer d'un volume de données cohérent et réaliste pour valider
l'ensemble du pipeline ETL et les calculs de KPI.

\subsubsection{Mécanisme d'extraction}

Les données sont extraites depuis les API REST de chaque application source (OPM, PTE, PMA)
via le client HTTP \texttt{httpx}. Chaque document brut est stocké dans MongoDB avec :
\begin{itemize}[leftmargin=1.5cm]
  \item Un identifiant de run ETL (\texttt{etl\_run\_id}) pour la traçabilité
  \item Un horodatage d'ingestion (\texttt{ingested\_at})
  \item Une clé d'unicité \texttt{(source\_id, etl\_run\_id)} pour le replay sans perte
\end{itemize}

Les collections Bronze suivent la convention \texttt{raw\_\{source\}\_\{entité\}}
(ex. \texttt{raw\_opm\_tickets}, \texttt{raw\_pte\_leaves}, \texttt{raw\_pma\_tasks}).

Pour les entités à pagination incomplète (congés PTE, réclamations PMA), le pipeline
effectue une lecture directe sur le MongoDB source en complément de l'API REST.

\subsection{Couche Argent — Transformation}

Chaque document brut est normalisé :
\begin{itemize}[leftmargin=1.5cm]
  \item Parsing des dates (ISO 8601, formats mixtes OPM/PTE/PMA)
  \item Résolution des identifiants imbriqués (dénormalisation des objets référencés)
  \item Validation des champs obligatoires, suppression des doublons par \texttt{source\_id}
  \item Chargement par \textit{upsert} en MongoDB (collections \texttt{stg\_*})
\end{itemize}

\subsection{Couche Or — Chargement Dimensionnel}

À partir des données Argent, le pipeline alimente :
\begin{enumerate}[leftmargin=1.5cm]
  \item \textbf{MDM} (Master Data Management) : déduplication cross-app des personnes
        via l'e-mail normalisé → une seule ligne \texttt{dim\_person} par individu
  \item \textbf{Dimensions} : chargement par upsert dans les tables \texttt{core.*}
  \item \textbf{Tables de faits} : insertion par batch (COPY ou INSERT ON CONFLICT IGNORE)
  \item \textbf{Data Marts} : \texttt{REFRESH MATERIALIZED VIEW} des 25 vues KPI
\end{enumerate}

% ═════════════════════════════════════════════════════════════════════
\section{Modèle Dimensionnel — Constellation}

\subsection{Choix du modèle}

Le modèle \textbf{Constellation} (ou \textit{Galaxy Schema}) est une généralisation du
modèle en étoile de Kimball. Il est composé de \textbf{plusieurs tables de faits}
partageant des \textbf{dimensions conformées} (conformed dimensions).

Ce choix est justifié par la nature multi-source du projet :
\begin{itemize}[leftmargin=1.5cm]
  \item Trois systèmes sources indépendants (OPM, PMA, PTE) génèrent trois faits distincts
  \item Les employés existent dans les trois systèmes → \texttt{dim\_person} est une
        \textbf{dimension conformée} (MDM unifié par e-mail normalisé)
  \item Le temps (\texttt{dim\_date}) est partagé par toutes les tables de faits
  \item Les marts exécutifs font des jointures \textbf{cross-facts} via \texttt{dim\_person}
\end{itemize}

\subsection{Schéma du modèle Constellation}

\vspace{0.5em}
\begin{center}
\begin{tikzpicture}[node distance=0.6cm and 1.8cm, >=Latex, font=\footnotesize]

  %% ── Fact tables (centre) ─────────────────────────────────────────
  \node[factbox, fill=opmcolor!10, draw=opmcolor] (fopm)
    {fact\_opm\_ticket\\{\tiny OPM}};

  \node[factbox, fill=pmacolor!10, draw=pmacolor, below=2.2cm of fopm] (fpma)
    {fact\_pma\_task\\{\tiny PMA}};

  \node[factbox, fill=ptecolor!10, draw=ptecolor, below=2.2cm of fpma] (fpte)
    {fact\_pte\_event\\{\tiny PTE}};

  %% ── Conformed dimensions (right) ─────────────────────────────────
  \node[confdim, right=3.5cm of fpma] (dperson)
    {\color{ptecolor}dim\_person\\{\tiny (MDM conforme)}\\{\tiny + dim\_person\_xref}};

  \node[confdim, above=0.5cm of dperson] (ddate)
    {\color{ptecolor}dim\_date\\{\tiny (temps conforme)}};

  \node[confdim, below=0.5cm of dperson] (ddept)
    {\color{ptecolor}dim\_department};

  \node[confdim, below=0.5cm of ddept] (dstatus)
    {\color{ptecolor}dim\_status\\{\tiny (junk dim)}};

  %% ── OPM-specific dimensions (left) ───────────────────────────────
  \node[dimbox, left=2.8cm of fopm, yshift=1.8cm] (dcontract)
    {dim\_contract};

  \node[dimbox, left=2.8cm of fopm] (dsite)
    {dim\_site};

  \node[dimbox, left=2.8cm of fopm, yshift=-1.8cm] (dequip)
    {dim\_equipment};

  \node[dimbox, above=0.4cm of dcontract] (dsource)
    {dim\_source\_system};

  \node[dimbox, above=0.4cm of dsource] (dtod)
    {dim\_time\_of\_day};

  %% ── PMA-specific dimensions (left) ───────────────────────────────
  \node[dimbox, left=2.8cm of fpma, yshift=0.5cm] (dproject)
    {dim\_project};

  \node[dimbox, left=2.8cm of fpma, yshift=-0.8cm] (dpriority)
    {dim\_priority};

  %% ── PTE-specific dimensions (left) ───────────────────────────────
  \node[dimbox, left=2.8cm of fpte, yshift=0.5cm] (dvehicle)
    {dim\_vehicle};

  \node[dimbox, left=2.8cm of fpte] (droom)
    {dim\_room};

  \node[dimbox, left=2.8cm of fpte, yshift=-0.8cm] (dleavetype)
    {dim\_leave\_type};

  %% ── Arrows: OPM fact → dims ──────────────────────────────────────
  \draw[arrow] (fopm) -- (dcontract);
  \draw[arrow] (fopm) -- (dsite);
  \draw[arrow] (fopm) -- (dequip);
  \draw[arrow] (fopm) -- (dsource);
  \draw[arrow] (fopm) -- (dtod);

  %% ── Arrows: PMA fact → dims ──────────────────────────────────────
  \draw[arrow] (fpma) -- (dproject);
  \draw[arrow] (fpma) -- (dpriority);

  %% ── Arrows: PTE fact → dims ──────────────────────────────────────
  \draw[arrow] (fpte) -- (dvehicle);
  \draw[arrow] (fpte) -- (droom);
  \draw[arrow] (fpte) -- (dleavetype);

  %% ── Arrows to conformed dims ─────────────────────────────────────
  \draw[confarrow] (fopm) -- (ddate);
  \draw[confarrow] (fpma) -- (ddate);
  \draw[confarrow] (fpte) -- (ddate);

  \draw[confarrow] (fopm) -- (dperson);
  \draw[confarrow] (fpma) -- (dperson);
  \draw[confarrow] (fpte) -- (dperson);

  \draw[confarrow] (fpma) -- (ddept);
  \draw[confarrow] (fpte) -- (ddept);

  \draw[confarrow] (fopm) -- (dstatus);
  \draw[confarrow] (fpma) -- (dstatus);
  \draw[confarrow] (fpte) -- (dstatus);

  %% ── Legend ───────────────────────────────────────────────────────
  \node[below=0.5cm of fpte, xshift=-1cm, font=\tiny] (leg1)
    {\tikz\draw[arrow] (0,0) -- (0.6,0); dim. spécifique};
  \node[right=0.3cm of leg1, font=\tiny]
    {\tikz\draw[confarrow] (0,0) -- (0.6,0); dim. conformée (partagée)};

\end{tikzpicture}
\end{center}

% ═════════════════════════════════════════════════════════════════════
\section{Dimensions}

\subsection{Dimensions conformées (partagées par tous les faits)}

\begin{tabularx}{\textwidth}{>{\bfseries}p{3.5cm} X}
\toprule
Dimension & Description \\
\midrule
\texttt{dim\_date}         & Calendrier complet (jour, semaine, mois, trimestre, année, week-end).
                             Clé entière \texttt{YYYYMMDD}. \\
\texttt{dim\_time\_of\_day}& Granularité horaire (0--23h). \\
\texttt{dim\_person}       & \textbf{Dimension MDM} : une ligne par individu réel, dédupliqué
                             par e-mail normalisé entre OPM, PMA et PTE. Contient : nom complet,
                             département, titre, ancienneté, statut actif. \\
\texttt{dim\_person\_xref} & Table de correspondance \texttt{(système source, ID source)} →
                             \texttt{person\_sk}. Méthode de correspondance : \textit{email}. \\
\texttt{dim\_department}   & Référentiel des départements (nom, centre de coût). \\
\texttt{dim\_source\_system}& Référentiel des systèmes sources (OPM, PMA, PTE). \\
\texttt{dim\_status}       & \textbf{Dimension poubelle (Junk Dim)} : combinaisons de flags
                             booléens et de statuts faible cardinalité :
                             \texttt{ticket\_status}, \texttt{task\_status}, \texttt{event\_status},
                             \texttt{is\_overdue}, \texttt{is\_sla\_breach}, \texttt{is\_accepted}. \\
\bottomrule
\end{tabularx}

\subsection{Dimensions spécifiques OPM}

\begin{tabularx}{\textwidth}{>{\bfseries}p{3.5cm} X}
\toprule
Dimension & Attributs principaux \\
\midrule
\texttt{dim\_contract} & Numéro de contrat, type, nature, SLA (heures), dates,
                         nom client, pays client, commercial (dénormalisé). \\
\texttt{dim\_site}     & Nom, adresse, coordonnées GPS (lat/lon), nom client (dénormalisé). \\
\texttt{dim\_equipment}& Numéro de série, nom, type (HARD/SOFT), version, constructeur,
                         site et contrat (dénormalisés). \\
\bottomrule
\end{tabularx}

\subsection{Dimensions spécifiques PMA}

\begin{tabularx}{\textwidth}{>{\bfseries}p{3.5cm} X}
\toprule
Dimension & Attributs principaux \\
\midrule
\texttt{dim\_project}  & Nom, type, priorité, chef de projet, client, dates début/fin,
                         date de clôture, compteur de réclamations. \\
\texttt{dim\_priority} & Niveau de priorité (Haute, Moyenne, Basse) + poids pour le scoring. \\
\bottomrule
\end{tabularx}

\subsection{Dimensions spécifiques PTE}

\begin{tabularx}{\textwidth}{>{\bfseries}p{3.5cm} X}
\toprule
Dimension & Attributs principaux \\
\midrule
\texttt{dim\_vehicle}    & Immatriculation, modèle, type de véhicule. \\
\texttt{dim\_room}       & Libellé, localisation, capacité (places). \\
\texttt{dim\_leave\_type}& Code congé, libellé, indicateur congé payé. \\
\bottomrule
\end{tabularx}

% ═════════════════════════════════════════════════════════════════════
\section{Tables de Faits}

\subsection{\texttt{fact\_opm\_ticket} — Source OPM}

\textit{Grain : un enregistrement par ticket d'intervention.}

\begin{tabularx}{\textwidth}{>{\ttfamily\footnotesize}p{4.5cm} >{\footnotesize}X}
\toprule
\textbf{Colonne} & \textbf{Description} \\
\midrule
source\_id              & Identifiant MongoDB du ticket source \\
ticket\_number          & Référence métier du ticket \\
created\_date\_sk       & FK → \texttt{dim\_date} (date de création) \\
assigned\_date\_sk      & FK → \texttt{dim\_date} (date d'assignation) \\
resolved\_date\_sk      & FK → \texttt{dim\_date} (date de résolution) \\
closed\_date\_sk        & FK → \texttt{dim\_date} (date de clôture) \\
created\_hour\_sk       & FK → \texttt{dim\_time\_of\_day} \\
contract\_sk            & FK → \texttt{dim\_contract} \\
site\_sk                & FK → \texttt{dim\_site} \\
equipment\_sk           & FK → \texttt{dim\_equipment} \\
client\_person\_sk      & FK → \texttt{dim\_person} (client demandeur) \\
assigned\_tech\_person\_sk & FK → \texttt{dim\_person} (technicien assigné) \\
status\_sk              & FK → \texttt{dim\_status} \\
time\_to\_assign\_min   & Mesure : délai d'assignation (minutes) \\
time\_to\_resolve\_min  & Mesure : délai de résolution (minutes) \\
time\_to\_close\_min    & Mesure : délai de clôture (minutes) \\
reopen\_count           & Mesure : nombre de réouvertures \\
\bottomrule
\end{tabularx}

\subsection{\texttt{fact\_pma\_task} — Source PMA}

\textit{Grain : un enregistrement par tâche de projet.}

\begin{tabularx}{\textwidth}{>{\ttfamily\footnotesize}p{4.5cm} >{\footnotesize}X}
\toprule
\textbf{Colonne} & \textbf{Description} \\
\midrule
source\_id              & Identifiant MongoDB de la tâche source \\
task\_ref               & Référence métier de la tâche \\
project\_sk             & FK → \texttt{dim\_project} \\
executor\_person\_sk    & FK → \texttt{dim\_person} (ingénieur exécutant) \\
team\_leader\_person\_sk& FK → \texttt{dim\_person} (chef de projet) \\
department\_sk          & FK → \texttt{dim\_department} \\
start\_date\_sk         & FK → \texttt{dim\_date} (début de tâche) \\
deadline\_date\_sk      & FK → \texttt{dim\_date} (échéance) \\
closed\_date\_sk        & FK → \texttt{dim\_date} (date de clôture) \\
status\_sk              & FK → \texttt{dim\_status} \\
priority\_sk            & FK → \texttt{dim\_priority} \\
progress\_pct           & Mesure : avancement en pourcentage \\
note                    & Mesure : note de qualité de la tâche \\
rating\_weight          & Mesure : poids pour le calcul du score PM \\
duration\_days          & Mesure : durée planifiée (jours) \\
\bottomrule
\end{tabularx}

\subsection{\texttt{fact\_pte\_event} — Source PTE}

\textit{Grain : un enregistrement par événement RH ou terrain.}\\
\textit{Types d'événements : \texttt{leave}, \texttt{vehicle\_usage}, \texttt{room\_reservation},
\texttt{vm\_request}, \texttt{intervention}.}

\begin{tabularx}{\textwidth}{>{\ttfamily\footnotesize}p{4.5cm} >{\footnotesize}X}
\toprule
\textbf{Colonne} & \textbf{Description} \\
\midrule
source\_id              & Identifiant MongoDB de l'événement source \\
event\_type             & Type d'événement (congé, véhicule, salle, VM, intervention) \\
event\_date\_sk         & FK → \texttt{dim\_date} (date de début) \\
end\_date\_sk           & FK → \texttt{dim\_date} (date de fin) \\
event\_hour\_sk         & FK → \texttt{dim\_time\_of\_day} \\
applicant\_person\_sk   & FK → \texttt{dim\_person} (demandeur/employé) \\
engineer\_person\_sk    & FK → \texttt{dim\_person} (ingénieur terrain) \\
department\_sk          & FK → \texttt{dim\_department} \\
vehicle\_sk             & FK → \texttt{dim\_vehicle} (si événement véhicule) \\
room\_sk                & FK → \texttt{dim\_room} (si réservation salle) \\
leave\_type\_sk         & FK → \texttt{dim\_leave\_type} (si congé) \\
status\_sk              & FK → \texttt{dim\_status} \\
duration\_min           & Mesure : durée (minutes) \\
km                      & Mesure : kilométrage (véhicule) \\
ram\_gb / disk\_gb      & Mesure : ressources VM allouées \\
business\_days          & Mesure : nombre de jours ouvrables (congé) \\
\bottomrule
\end{tabularx}

% ═════════════════════════════════════════════════════════════════════
\section{Data Marts — Vues Matérialisées}

Les data marts sont des \textbf{vues matérialisées PostgreSQL} pré-calculées,
regroupées en 4 schémas. Elles sont rafraîchies à chaque fin de run ETL.

\subsection{mart\_opm — 7 vues}

\begin{tabularx}{\textwidth}{>{\ttfamily\footnotesize}p{5cm} >{\footnotesize}X}
\toprule
\textbf{Vue} & \textbf{Calcul} \\
\midrule
v\_tickets\_by\_status     & COUNT tickets GROUP BY statut et période \\
v\_sla\_compliance         & \% tickets sans dépassement SLA par contrat et période \\
v\_mttr                    & Moyenne(\texttt{time\_to\_resolve\_min} / 60) par contrat \\
v\_mtta                    & Moyenne(\texttt{time\_to\_assign\_min} / 60) par contrat \\
v\_first\_call\_resolution & \% tickets fermés sans réouverture (\texttt{reopen\_count = 0}) \\
v\_technician\_load        & COUNT tickets par technicien (JOIN \texttt{dim\_person}) \\
v\_contract\_health        & \% conformité SLA par contrat (identique à SLA, usage distinct) \\
\bottomrule
\end{tabularx}

\subsection{mart\_pma — 7 vues}

\begin{tabularx}{\textwidth}{>{\ttfamily\footnotesize}p{5cm} >{\footnotesize}X}
\toprule
\textbf{Vue} & \textbf{Calcul} \\
\midrule
v\_portfolio\_status        & COUNT projets distincts par statut de tâche et période \\
v\_tasks\_by\_status        & COUNT tâches GROUP BY statut \\
v\_tasks\_by\_priority      & COUNT tâches GROUP BY priorité (JOIN \texttt{dim\_priority}) \\
v\_on\_time\_delivery       & \% tâches clôturées sans retard (\texttt{is\_overdue = false}) \\
v\_team\_leader\_score      & Moyenne pondérée(\texttt{note × rating\_weight}) par chef de projet \\
v\_engineer\_productivity   & COUNT tâches par ingénieur et période \\
v\_overdue\_index           & \% tâches en retard par projet \\
\bottomrule
\end{tabularx}

\subsection{mart\_pte — 6 vues}

\begin{tabularx}{\textwidth}{>{\ttfamily\footnotesize}p{5cm} >{\footnotesize}X}
\toprule
\textbf{Vue} & \textbf{Calcul} \\
\midrule
v\_leave\_consumption       & SUM(\texttt{business\_days}) par département, type congé, période \\
v\_headcount\_active        & COUNT employés actifs internes par département \\
v\_vehicle\_utilization     & SUM(\texttt{km}) par immatriculation et période \\
v\_room\_occupancy          & SUM(\texttt{duration\_min} / 60) heures d'occupation par salle \\
v\_vm\_lead\_time           & Moyenne(\texttt{duration\_min} / 60) pour les requêtes VM \\
v\_intervention\_throughput & COUNT interventions par ingénieur et période \\
\bottomrule
\end{tabularx}

\subsection{mart\_exec — 5 vues transversales}

Ces vues agrègent des données \textbf{cross-sources} via la dimension conformée
\texttt{dim\_person}.

\begin{tabularx}{\textwidth}{>{\ttfamily\footnotesize}p{5cm} >{\footnotesize}X}
\toprule
\textbf{Vue} & \textbf{Calcul} \\
\midrule
v\_company\_throughput  & UNION tickets résolus (OPM) + tâches clôturées (PMA) par période \\
v\_pm\_scorecard        & JOIN \texttt{fact\_pma\_task} ↔ \texttt{fact\_opm\_ticket} via
                          \texttt{dim\_person} (chef de projet → ingénieurs → tickets) ;
                          calcule livraison à temps + délai moyen de résolution \\
v\_revenue\_at\_risk    & Contrats expirant dans 90 jours ayant eu une rupture SLA
                          au trimestre précédent \\
v\_client\_health       & \% conformité SLA par client (nom dénormalisé sur \texttt{dim\_contract}) \\
v\_workforce\_availability & COUNT employés actifs ayant un événement PTE par département \\
\bottomrule
\end{tabularx}

% ═════════════════════════════════════════════════════════════════════
\section{Récapitulatif Global}

\begin{tabularx}{\textwidth}{>{\bfseries}p{5cm} >{\centering\arraybackslash}p{1.8cm} X}
\toprule
\textbf{Objet} & \textbf{Nombre} & \textbf{Localisation} \\
\midrule
Schémas PostgreSQL         & 5  & \texttt{core}, \texttt{mart\_opm}, \texttt{mart\_pma}, \texttt{mart\_pte}, \texttt{mart\_exec} \\
Dimensions conformées      & 7  & \texttt{core.*} \\
Dimensions spécifiques     & 8  & \texttt{core.*} \\
Tables de faits            & 3  & \texttt{core.fact\_opm\_ticket}, \texttt{core.fact\_pma\_task}, \texttt{core.fact\_pte\_event} \\
Vues matérialisées (KPI)   & 25 & \texttt{mart\_opm} (7) + \texttt{mart\_pma} (7) + \texttt{mart\_pte} (6) + \texttt{mart\_exec} (5) \\
Collections MongoDB Bronze & 15 & \texttt{raw\_*} \\
Collections MongoDB Argent & 15 & \texttt{stg\_*} \\
\bottomrule
\end{tabularx}

\end{document}
